% !TEX program = xelatex
\documentclass[UTF8,11pt,a4paper,fontset=none]{ctexart}

\usepackage[margin=2.15cm,top=2.1cm,bottom=2.2cm]{geometry}
\usepackage{fontspec}
\usepackage{xeCJK}
\usepackage{amsmath,amssymb}
\usepackage{booktabs}
\usepackage{tabularx}
\usepackage{longtable}
\usepackage{array}
\usepackage{enumitem}
\usepackage{xcolor}
\usepackage{hyperref}
\usepackage{fancyhdr}
\usepackage{titlesec}
\usepackage[most]{tcolorbox}
\usepackage{listings}
\usepackage{microtype}
\usepackage{ragged2e}

\setmainfont[
  Path=/usr/local/texlive/2024/texmf-dist/fonts/opentype/public/tex-gyre/,
  Extension=.otf,
  UprightFont=*-regular,
  BoldFont=*-bold,
  ItalicFont=*-italic,
  BoldItalicFont=*-bolditalic
]{texgyretermes}
\setsansfont[
  Path=/usr/local/texlive/2024/texmf-dist/fonts/opentype/public/tex-gyre/,
  Extension=.otf,
  UprightFont=*-regular,
  BoldFont=*-bold,
  ItalicFont=*-italic,
  BoldItalicFont=*-bolditalic
]{texgyreheros}
\setmonofont[
  Path=/usr/local/texlive/2024/texmf-dist/fonts/opentype/public/tex-gyre/,
  Extension=.otf,
  UprightFont=*-regular,
  BoldFont=*-bold,
  ItalicFont=*-italic,
  BoldItalicFont=*-bolditalic
]{texgyrecursor}
\setCJKmainfont[Path=/System/Library/Fonts/]{Hiragino Sans GB.ttc}
\setCJKsansfont[Path=/System/Library/Fonts/]{Hiragino Sans GB.ttc}
\setCJKmonofont[Path=/Library/Fonts/]{Arial Unicode.ttf}

\definecolor{navy}{HTML}{17324D}
\definecolor{steel}{HTML}{35627A}
\definecolor{lightsteel}{HTML}{EEF5F7}
\definecolor{softgray}{HTML}{F7F8FA}
\definecolor{darkgray}{HTML}{444444}
\definecolor{riskred}{HTML}{9B2D30}
\definecolor{goodgreen}{HTML}{1F6B4A}

\hypersetup{
  colorlinks=true,
  linkcolor=navy,
  urlcolor=steel,
  citecolor=steel,
  pdfauthor={OpenAI},
  pdftitle={Pricing \& Promotion Decision Engine 项目计划},
  pdfsubject={定价与促销决策 Demo 计划}
}

\pagestyle{fancy}
\setlength{\headheight}{14pt}
\fancyhf{}
\lhead{Pricing \& Promotion Decision Engine}
\rhead{项目执行计划}
\cfoot{\thepage}
\renewcommand{\headrulewidth}{0.4pt}
\renewcommand{\footrulewidth}{0pt}

\titleformat{\section}{\Large\bfseries\color{navy}}{\thesection}{0.75em}{}
\titleformat{\subsection}{\large\bfseries\color{steel}}{\thesubsection}{0.65em}{}
\titleformat{\subsubsection}{\normalsize\bfseries\color{darkgray}}{\thesubsubsection}{0.55em}{}
\titlespacing*{\section}{0pt}{1.0em}{0.55em}
\titlespacing*{\subsection}{0pt}{0.85em}{0.35em}
\titlespacing*{\subsubsection}{0pt}{0.65em}{0.25em}

\setlist[itemize]{leftmargin=1.5em,itemsep=0.25em,topsep=0.25em}
\setlist[enumerate]{leftmargin=1.65em,itemsep=0.25em,topsep=0.25em}
\renewcommand{\arraystretch}{1.22}
\newcolumntype{Y}{>{\RaggedRight\arraybackslash}X}
\newcolumntype{L}[1]{>{\RaggedRight\arraybackslash}p{#1}}
\newcolumntype{C}[1]{>{\centering\arraybackslash}p{#1}}

\tcbset{
  enhanced,
  sharp corners,
  boxrule=0.6pt,
  colback=softgray,
  colframe=steel,
  left=7pt,
  right=7pt,
  top=6pt,
  bottom=6pt,
  before skip=0.65em,
  after skip=0.65em
}
\newtcolorbox{decisionbox}[1]{colback=lightsteel,colframe=navy,title={#1},fonttitle=\bfseries}
\newtcolorbox{riskbox}[1]{colback=white,colframe=riskred,title={#1},fonttitle=\bfseries\color{riskred}}
\newtcolorbox{actionbox}[1]{colback=white,colframe=goodgreen,title={#1},fonttitle=\bfseries\color{goodgreen}}

\lstset{
  basicstyle=\ttfamily\small,
  breaklines=true,
  columns=fullflexible,
  frame=single,
  rulecolor=\color{steel},
  backgroundcolor=\color{softgray},
  xleftmargin=0.5em,
  xrightmargin=0.5em
}

\setlength{\parindent}{0pt}
\setlength{\parskip}{0.45em}
\emergencystretch=3em

\begin{document}

\begin{titlepage}
\centering
\vspace*{1.8cm}
{\Huge\bfseries\color{navy} Pricing \& Promotion\\Decision Engine\par}
\vspace{0.6cm}
{\Large 定价与促销决策 Demo 完整执行计划\par}
\vspace{1.1cm}
\begin{tcolorbox}[width=0.86\textwidth,colback=lightsteel,colframe=navy]
\large
核心叙事：给定有效价格、促销状态、成本、库存和需求弹性，系统估计需求变化，模拟反事实，并生成需要实验验证的受约束 raise-and-test 定价候选。
\end{tcolorbox}
\vspace{0.8cm}
\begin{tabular}{ll}
\toprule
项目定位 & 求职用 Decision Quant / Applied Scientist / Economist demo \\
主交付物 & Streamlit web app + GitHub repo + methodology + executive memo \\
方法主线 & 需求估计 + 反事实模拟 + 利润优化 + A/B 验证设计 \\
当前状态 & 数据清洗、EDA、需求估计、反事实优化、实验设计和 Streamlit MVP 已完成 \\
技术栈 & Python, pandas, statsmodels, scipy, Streamlit \\
\bottomrule
\end{tabular}
\vfill
{\large 版本：v1.1\quad 日期：\today\par}
\end{titlepage}

\setcounter{tocdepth}{1}
{\small\setlength{\parskip}{0pt}\tableofcontents}
\newpage

\section{一句话结论}

\begin{decisionbox}{最终方向}
不要做单纯销量预测 demo。做一个 \textbf{Pricing \& Promotion Decision Engine}：它估计价格和促销对需求的影响，模拟有效价格与促销状态反事实，并在成本、库存、最低销量、最低毛利等业务约束下生成可审计的 raise-and-test 定价候选。
\end{decisionbox}

这比“我预测销量”强，因为它直接回答招聘方关心的问题：\textbf{这个模型如何改变业务决策、候选动作的收益和风险是什么、上线前需要怎样验证？}

这个项目不把“发明一个新的需求模型”作为卖点。它的定位是把 scanner-data demand estimation、pricing/revenue optimization、promotion response 和 experimentation 这些已有研究脉络，工程化成一个可交互、可复现、能产生业务建议的 app。

项目的最终目标不是展示模型复杂度，而是展示四种能力：
\begin{enumerate}
  \item 会把业务问题转成可估计、可优化的决策问题；
  \item 会处理自身价格弹性、交叉价格弹性、促销条件效应、价格和促销内生性、库存约束等真实问题；
  \item 会把模型封装成业务方能使用的 web app；
  \item 会把模型建议转成实验方案，而不是把观察性回归当成最终因果结论。
\end{enumerate}

\subsection{当前项目状态}

截至 v1.1，项目已经从计划阶段推进到可运行 MVP。当前 repo 已包含：

\begin{itemize}
  \item \texttt{01\_data\_cleaning.ipynb}、\texttt{02\_eda.ipynb}、\texttt{03\_demand\_estimation.ipynb}、\texttt{04\_counterfactual.ipynb} 和 \texttt{05\_ab\_testing\_design.ipynb}；
  \item SKU-store-week panel、brand-size-store-week panel、demand modeling dataset、model coefficients、top recommendations 和 experiment candidates；
  \item Streamlit app，包括 Demand Model、Counterfactual Simulator、Profit Optimizer、Experiment Design、Limitations 和 Upload \& Score 六个页面；
  \item methodology Markdown / TeX / PDF，以及 data cleaning、EDA、demand model、counterfactual 和 A/B validation 报告。
\end{itemize}

下一步不再是“开始建 app”，而是收口：先做 \texttt{07\_cannibalization\_robustness.ipynb} 检查同品牌跨包装规格互食，再补 tests、case study 和部署；若时间充足，再做 \texttt{08\_causal\_iv.ipynb}。

\section{为什么不是纯预测，也不是选址 demo}

\subsection{Counterarguments}

\begin{enumerate}
  \item \textbf{纯预测回答的是“会卖多少”，不是“应该怎么做”。}公司招聘定价、增长、市场设计、运营科学岗位时，核心关心的是价格、折扣、补贴、排序、库存分配等决策变量。
  \item \textbf{选址是低频决策，反馈慢，实验难。}电商和平台公司更常见的高频决策是价格、促销、推荐、广告预算和供需匹配。选址 demo 更像城市规划或零售地产分析，求职信号较窄。
  \item \textbf{黑箱预测没有经济学优势。}如果主线只是 XGBoost 预测销量，你与普通机器学习候选人没有明显差异。你的优势是因果识别、反事实推理、均衡直觉和约束优化。
  \item \textbf{没有优化器的 app 只是 dashboard。}招聘方需要看到“输入参数 - 反事实 - 推荐决策 - 风险验证”的闭环。
\end{enumerate}

\subsection{项目判别标准}

\begin{tabularx}{\textwidth}{L{4.0cm}Y Y}
\toprule
项目类型 & 能回答的问题 & 求职信号 \\
\midrule
销量预测 & 下周销量是多少？ & 中等，偏 analyst/ML exercise \\
定价决策引擎 & 价格提高 5\% 后，销量、收入、利润怎样变？哪个价格值得测试？ & 高，直接对齐 decision quant \\
促销响应分析 & 打折是否真的增加利润？需要怎样实验验证？ & 高，对齐 experimentation / causal inference \\
选址分析 & 新店应开在哪里？ & 小众，低频，难验证 \\
\bottomrule
\end{tabularx}

\section{目标岗位与目标公司类型}

\subsection{岗位定位}

这个 demo 对齐以下岗位，而不是泛泛的数据分析岗位：

\begin{itemize}
  \item Applied Scientist, Pricing / Marketplace / Optimization
  \item Economist, Causal Inference / Experimentation / Pricing
  \item Data Scientist, Pricing / Growth / Monetization
  \item Decision Scientist / Revenue Optimization Analyst
  \item Product Data Scientist, Marketplace / Promotion / Customer Growth
\end{itemize}

\subsection{公司类型}

\begin{tabularx}{\textwidth}{L{3.2cm}Y Y}
\toprule
公司类型 & 典型业务杠杆 & Demo 对齐方式 \\
\midrule
电商/零售平台 & 价格、折扣、库存、推荐、促销预算 & SKU-store-week 需求估计 + 利润优化 \\
本地服务/餐饮 SaaS & 商家定价、促销、用户增长、订阅留存 & 促销响应 + A/B test 设计 \\
广告/营销科技 & 预算分配、投放 ROI、客户分群 & uplift 建模 + constrained optimization \\
平台经济公司 & 补贴、匹配、动态定价、供需平衡 & 反事实模拟 + 约束优化 \\
经济咨询 & 损害评估、反事实、政策模拟 & 识别问题 + 透明模型 + memo 写作 \\
\bottomrule
\end{tabularx}

\section{业务问题定义}

\subsection{核心问题}

该 demo 要模拟一个真实定价团队面对的问题：

\begin{quote}
对于给定商品或品类，当前价格和促销策略是否合理？如果改变有效价格或促销状态，需求、收入、利润、库存风险会如何变化？在成本和库存约束下，哪些价格和促销状态值得作为 raise-and-test 候选？上线前需要怎样设计实验验证？
\end{quote}

\subsection{必须回答的业务问题}

\begin{enumerate}
  \item 当前价格是否过高或过低？
  \item 促销提高的是销量、收入，还是利润？
  \item 在给定 unit cost 下，候选价格是否仍有正 margin？
  \item 在库存有限时，推荐策略是否造成 stockout？
  \item 若价格上调或下调 5\%，收入和利润分别如何变化？
  \item 候选价格是否需要 A/B test 验证？实验需要多少样本？
\end{enumerate}

\section{最小可行产品：MVP 架构}

\subsection{六个核心模块}

\begin{longtable}{L{3.2cm}L{4.0cm}L{7.6cm}}
\toprule
模块 & 输入 & 输出 \\
\midrule
\endfirsthead
\toprule
模块 & 输入 & 输出 \\
\midrule
\endhead
1. 数据处理 & 交易数据、商品信息、日期、促销标记 & SKU $\times$ store $\times$ week panel \\
2. 需求估计 & own price, competitor price, quantity, promo, SKU-store FE, week FE & 自身价格弹性、交叉价格弹性、促销条件效应、预测误差 \\
3. 反事实模拟 & 新有效价格、促销状态、成本、库存、需求冲击 & 预测销量、收入、利润、变化率 \\
4. 优化器 & 目标函数、约束、价格范围、促销状态 & raise-and-test 价格候选、促销状态、预期利润提升 \\
5. 实验设计 & baseline 指标、MDE、power、方差 & 样本量、实验周期、guardrail metrics \\
6. 应用层 & app 选择器、场景输入、用户上传数据 & 交互式结果页、上传数据评分、风险提示 \\
\bottomrule
\end{longtable}

\subsection{MVP 不做什么}

\begin{riskbox}{必须排除的范围}
\begin{itemize}
  \item 不做复杂全栈前端。Streamlit 足够。
  \item 不把深度学习作为主模型。主模型必须透明。
  \item 不做“只有图表没有决策建议”的 dashboard。
  \item 不声称观察性价格弹性就是因果效应。
  \item 不伪造真实成本和库存；如果数据没有这些字段，就作为 scenario inputs。
\end{itemize}
\end{riskbox}

\section{数据方案}

\subsection{数据来源：Dominick's Finer Foods, Cereals 品类}

本项目选用 \textbf{Dominick's Finer Foods (DFF)} 零售面板，范围锁定在 \textbf{Ready-to-Eat Cereals}（即食麦片）品类，覆盖全部门店与全部可用周次。数据由芝加哥大学 Booth 商学院 James M.\ Kilts Center 公开提供，时间范围约 1989--1997 年，包含约 85 家 Chicago 地区 Dominick's 门店的周级 UPC-store 交易记录。

\subsection{为什么选 Dominick's 而不是其他公开零售面板}

\begin{enumerate}
  \item \textbf{字段完整度最高。}DFF 同时含有 price、quantity (\texttt{MOVE})、promo code (\texttt{SALE})、profit margin (\texttt{PROFIT})、week、store、UPC。利润函数所需的 unit cost 可由 price 与 PROFIT 反推，不必作为场景假设。这是所有公开零售面板中唯一同时满足这四类字段的数据集。
  \item \textbf{学术与产业双重语境。}DFF 是实证定价、产业组织、营销科学领域使用最广的公开面板。向 economist / applied scientist / pricing data scientist 岗位的面试官提到 ``Dominick's'' 有即时识别度，不需要额外解释数据来源。
  \item \textbf{样本规模合适。}约 85 家门店 $\times$ 约 400 周 $\times$ 单品类 UPC，足够估计价格弹性、品牌间交叉弹性和促销条件效应，同时可以在本地笔记本上完成处理，不依赖云端。
  \item \textbf{促销字段精细。}\texttt{SALE} 代码区分非促销、bonus buy、coupon、simple price reduction 等形式，允许把纯价格变动效应与条件 sale-code effect 分离——多数公开数据集只有二元 promo flag。
\end{enumerate}

\subsection{为什么选 Cereals 品类}

\begin{tabularx}{\textwidth}{L{3.2cm}Y}
\toprule
维度 & Cereals 品类的优势 \\
\midrule
学术锚点 & Nevo (2001, \textit{Econometrica}) ``Measuring Market Power in the Ready-to-Eat Cereal Industry'' 是实证定价方向的 landmark paper；本项目使用 DFF Cereals 作为同一品类和 scanner-data 语境下的透明决策工程化版本，不声称复刻 Nevo 的结构模型或原始数据 \\
品牌结构 & Kellogg's、General Mills、Post、Quaker、Ralston、Nabisco、private label 等 5--7 家势均力敌，交叉价格弹性有经济含义 \\
利润叙事 & 品类 price-cost margin 约 40\%--45\%（Nevo 估计），是 market power 的典型案例，适合作为 memo 叙事支点 \\
UPC 归一化 & 包装形式以盒装为主，size 维度少，\texttt{price per oz} 归一化简单 \\
识别干扰 & 季节性弱、消费平稳、囤货 (stockpiling) 动机小，静态需求模型的跨期偏误较小 \\
促销密度 & 促销频率高，能提供充足的价格内变动，有利于弹性识别 \\
\bottomrule
\end{tabularx}

\subsection{为什么排除其他候选品类}

\begin{itemize}
  \item \textbf{Soft Drinks：}2L 瓶 / 12 罐装 / 6 罐装 / 20oz 瓶等包装形式复杂，UPC 归一化工作量大；夏季季节性强，时间固定效应吸收不干净；消费者对促销的囤货行为 (Hendel-Nevo 效应) 会导致静态需求模型系统性低估真实弹性。
  \item \textbf{Beer：}税收、节假日与州法规冲击大，识别更复杂；部分时期促销折扣形式受法规限制。
  \item \textbf{Analgesics / Bath Soap：}品牌忠诚度极高、弹性小、价格内变动少，不足以做有信息量的弹性估计。
  \item \textbf{Frozen / Canned Foods：}促销频率低，价格变动不充分。
  \item \textbf{跨品类合并：}不同品类弹性异质，合并估计单一 $\beta$ 在经济上错误；按品类分别估计则实质是多个独立项目，不如做深一个品类。
\end{itemize}

\subsection{数据基本信息}

\begin{tabularx}{\textwidth}{L{4.0cm}Y}
\toprule
维度 & 取值 \\
\midrule
来源 & UChicago Booth, James M.\ Kilts Center for Marketing --- Dominick's Finer Foods Database \\
公开链接 & \url{https://www.chicagobooth.edu/research/kilts/research-data/dominicks} \\
品类 & Ready-to-Eat Cereals（交易文件前缀 \texttt{wcer}，UPC 映射 \texttt{upccer}）\\
时间范围 & 约 1989--1997 年，周级 \\
门店数 & 约 85 家 Chicago 地区门店 \\
UPC 数量 & UPC 映射表 490 行；Step 2 清洗后交易面板保留 486 个有有效交易记录的 UPC，可继续聚合到 brand $\times$ size 主 SKU \\
主要品牌 & Kellogg's, General Mills, Post, Quaker, Ralston, Nabisco, private label \\
核心字段 & \texttt{UPC, STORE, WEEK, MOVE, PRICE, PROFIT, SALE, QTY} \\
可反推字段 & \texttt{unit\_price} = \texttt{PRICE}/\texttt{QTY}；\texttt{unit\_cost} = \texttt{unit\_price} $\times$ (1 - \texttt{PROFIT}/100) \\
规模估计 & Step 2 清洗后为 4,654,156 行，覆盖 486 UPC、93 stores、366 weeks、36,199 UPC-store pairs；每对中位覆盖 78 周 \\
\bottomrule
\end{tabularx}

\subsection{字段结构与 DFF 对应关系}

项目主面板的最小字段及其与 DFF 原始字段的映射如下：

\begin{tabularx}{\textwidth}{L{2.8cm}L{3.6cm}Y L{2.4cm}}
\toprule
项目字段 & DFF 对应 & 含义 & 备注 \\
\midrule
product\_id & \texttt{UPC} & 商品 ID & 可聚合到 brand \\
store\_id & \texttt{STORE} & 门店 ID & 必须，用于 SKU-store FE \\
week & \texttt{WEEK} & 周序号 & DFF 自定义周计数 \\
price & \texttt{PRICE}/\texttt{QTY} & 有效单位销售价格（美元） & 价格弹性主变量 \\
price\_per\_oz & \texttt{price}/\texttt{size\_oz} & 单位容量价格 & 用于跨品牌比较 \\
bundle\_price & \texttt{PRICE} & 捆绑销售价格（美元） & 用于审计原始价格 \\
quantity & \texttt{MOVE} & 周销售单位数 & 单位数而非重量 \\
promo & \texttt{SALE} & 促销代码 & MVP 转为二元 indicator \\
cost & \texttt{unit\_price} $\times$ (1 - \texttt{PROFIT}/100) & 单位成本 & 由 DFF gross margin 反推 \\
brand & 由 UPC 映射 & 品牌 & 用于异质性与交叉弹性 \\
size\_oz & 由 UPC 映射 & 包装容量（盎司） & 归一化 price per oz 用 \\
category & 固定为 \texttt{cer} & 品类 & 本项目锁定 Cereals \\
inventory & 数据缺 & 库存 & 作为 scenario input \\
competitor price & 同 store-week 内其他品牌均价 & 竞争价格 & 用基准份额加权 \\
\bottomrule
\end{tabularx}

\subsection{数据质量底线}

\begin{itemize}
  \item 不使用只有销量而缺失价格的周；
  \item 不使用只有浏览 / 曝光而无交易结果的数据；
  \item 不使用纯合成数据；DFF 为真实零售面板；
  \item \texttt{PROFIT} 字段缺失或异常的记录不进入利润计算；
  \item 库存字段 DFF 不提供，作为 scenario input 在 app 中明确标注；
  \item 所有聚合与过滤决策（UPC 过滤阈值、最小周数、异常价格规则）写入 \texttt{rawData/README.md} 或 \texttt{reports/data\_cleaning\_summary.md}，并在根目录 \texttt{README.md} 中链接，保证可复现。
\end{itemize}

\subsection{Step 2 数据诊断记录}

第一轮 notebook 清洗已经产出 \texttt{data/processed/sku\_store\_week\_panel.parquet}、\texttt{reports/data\_cleaning\_summary.md} 和 \texttt{reports/figures/data\_quality\_wcer.png}。进入模块化前，以下数据事实需要固定到文档和清洗规则中：

\begin{itemize}
  \item \texttt{PRICE=0} 与 \texttt{MOVE=0} 几乎完全重合：raw 中所有 1,850,703 个零销量行都同时为零价格；另有 677 行 \texttt{PRICE=0, MOVE>0}，视为无效价格记录。因此主需求样本继续使用 \texttt{PRICE>0} 和 \texttt{MOVE>0} 过滤。
  \item \texttt{SALE} 除 B/S/C 外还出现 \texttt{G} 和 \texttt{L}。其中 \texttt{G} 频次足够高，MVP 二元 \texttt{promo} 中计为促销，并在分类分析中保留为 \texttt{unknown\_promo}；\texttt{L} 只有 1 行，并入缺失。
  \item \texttt{PROFIT} 中位数为 13.7\%，validation sanity range 应使用 $[5,30]$，而不是 $[20,50]$；负利润或大于等于 99\% 的记录不进入成本和利润计算。
  \item \texttt{COM\_CODE} 对所有 Cereals UPC 都是 311，不存在需要额外处理的子品类混杂。
  \item \texttt{SIZE} 严格 OZ 正则可解析率为 72.2\%。\texttt{11.25O}、\texttt{1.25 O}、\texttt{19.25Z} 这类截断单位要用鲁棒正则处理；\texttt{ASST}、\texttt{1 CT}、\texttt{end} 等脏值要从 \texttt{price\_per\_oz} 和竞争价格构造中剔除或人工修复。
  \item 清洗后样本覆盖 486 UPC、93 stores、366 weeks 和 36,199 UPC-store pairs；每对中位覆盖 78 周，足以支持 SKU-store pair fixed effects。
\end{itemize}

\subsection{当前已完成的实证产出}

当前项目已经完成 Step 2--Step 5 的主要流水线，关键结果需要同步到 README、methodology、app 和报告中：

\begin{itemize}
  \item \textbf{EDA 与聚合：}\texttt{brand\_size\_store\_week\_panel.parquet} 已生成，brand-size 粒度作为 baseline，UPC 粒度作为 robustness；brand mapping 经 DESCRIP 规则与 manufacturer code 交叉验证。
  \item \textbf{需求估计：}baseline 估计得到 $\widehat{\beta}_{own}\approx -1.75$，within $R^2$ 约 0.75；这里的 within $R^2$ 指 fixed-effects panel regression fit，不是 out-of-sample predictive accuracy。加入竞争价格后 $\widehat{\beta}_{own}\approx -1.73$、$\widehat{\beta}_{cross}\approx 0.65$、促销 sale-code 系数约 $0.43$ log points，对应 $\exp(0.43)-1\approx54\%$ 的条件销量差异。UPC robustness 给出 $\widehat{\beta}_{own}\approx -1.90$、$\widehat{\beta}_{cross}\approx 0.50$。
  \item \textbf{预测与重转换：}log-demand 的 level prediction 使用 Duan smearing；cell-anchored counterfactual 从 observed level baseline 出发，因此不再额外乘 smearing factor，避免重复校正。
  \item \textbf{反事实优化：}eligible cells 中 98.5\% 的候选动作贴到上界 price guardrail。该结果必须解释为模型诊断和 raise-and-test candidates，而不是 production deployment recommendations。
  \item \textbf{实验设计：}top candidates 多数属于 high-risk / underpowered single-store flights，说明优化器候选动作必须通过更长周期或更大范围实验验证。
\end{itemize}

\subsection{处理流程}

\begin{enumerate}
  \item 从 Kilts Center 下载 \texttt{wcer.csv}（交易）与 \texttt{upccer.csv}（UPC 映射）；
  \item 清洗异常价格、零销量、重复 (UPC, store, week) 记录；主需求模型仅使用 $Q>0$ 的周；
  \item 构造有效单位价格 \texttt{unit\_price} = \texttt{PRICE}/\texttt{QTY}，并用它作为需求模型中的 price；
  \item 构造 \texttt{unit\_cost} = \texttt{unit\_price} $\times$ (1 - \texttt{PROFIT}/100)；
  \item 由 UPC 映射表补充 brand 与 size\_oz，归一化 \texttt{price\_per\_oz}；
  \item 按 brand $\times$ size\_oz 聚合（或直接 brand 级），降低 UPC 噪声；但聚合可能把同一 brand-size 内不同 UPC 的价格、促销和 bundle 组合变化混进需求反应，因此 UPC 粒度结果必须作为 robustness check；
  \item 构造 SKU-store-week panel，计算 log price、log quantity、二元 promo indicator；
  \item MVP 不单独模拟连续 discount depth；若增强版加入折扣深度，必须先显式估计对应系数；
  \item 构造同一 store-week 内其他品牌的加权 \texttt{price\_per\_oz} 作为竞争价格控制；权重使用 pre-period 或滚动滞后销量份额，不使用当期销量份额；
  \item 过滤样本过少的 SKU（如少于 20 周观察）；
  \item 按时间划分 train / holdout（如最后 20\% 周留作 holdout）；
  \item 保存 processed panel 为 parquet，供模型与 app 调用。
\end{enumerate}

\subsection{Brand-size 聚合公式}

当 UPC 粒度观测聚合到 brand $\times$ size\_oz $\times$ store $\times$ week 时，令 $g$ 表示 brand-size cell，$u\in g$ 表示属于该 cell 的 UPC。聚合规则必须固定，避免后续模型变量来源不清：

\begin{equation}
Q_{gst}=\sum_{u\in g} MOVE_{ust}.
\end{equation}

销售额先在 UPC 层根据 DFF bundle price 与 bundle quantity 计算，再聚合：

\begin{equation}
Sales_{gst}=\sum_{u\in g}\frac{PRICE_{ust}\times MOVE_{ust}}{QTY_{ust}}.
\end{equation}

聚合后的有效价格是 revenue-per-unit，也等价于按销量加权的有效单位价格：

\begin{equation}
P^{eff}_{gst}=\frac{Sales_{gst}}{Q_{gst}}.
\end{equation}

聚合后的 unit-cost proxy 按销量加权：

\begin{equation}
c_{gst}=
\frac{\sum_{u\in g} c_{ust} MOVE_{ust}}
{\sum_{u\in g} MOVE_{ust}}.
\end{equation}

baseline 促销变量使用 any-promo indicator：

\begin{equation}
Promo^{any}_{gst}=\mathbf{1}\left\{\sum_{u\in g} Promo_{ust}>0\right\},
\end{equation}

同时保留按销量加权的促销占比作为诊断变量：

\begin{equation}
Promo^{share}_{gst}
=
\frac{\sum_{u\in g} Promo_{ust}MOVE_{ust}}
{\sum_{u\in g} MOVE_{ust}}.
\end{equation}

\subsection{竞争价格构造}

竞争价格用同一门店、同一周内其他品牌的单位容量价格构造：

\begin{equation}
P^{comp}_{ist}=\sum_{b\ne B(i)} w^{base}_{bs}\,P^{oz}_{bst}.
\end{equation}

其中 $B(i)$ 是商品 $i$ 所属品牌，$P^{oz}_{bst}$ 是品牌 $b$ 在门店 $s$、周 $t$ 的加权平均 \texttt{price\_per\_oz}。权重 $w^{base}_{bs}$ 使用 pre-period 或滚动滞后销量份额，并在当前周 $t$ 之前确定；不使用当期销量份额，避免把当期需求冲击机械地放进竞争价格变量。若某门店历史销量不足，MVP 可退回到其他品牌等权平均。

\section{Relation to Prior Work}

这个项目的研究问题不是全新的。价格弹性、促销反应、品牌替代和利润优化都已经有成熟文献和商业系统。项目的价值在于：用公开 DFF 数据做一个透明、可运行、可解释的 decision engine，把研究问题从“估计需求”推进到“生成候选价格、解释风险、设计实验”。

\subsection{Scanner Data 与零售定价研究}

Dominick's Finer Foods 数据本身来自 UChicago Booth Kilts Center 的 store-level shelf management and pricing research 合作，包含多品类、门店级、周级 UPC movement data。它适合作为 MVP 数据源，不是因为它最新，而是因为字段足够接近真实定价团队需要的最小闭环：价格、销量、促销代码、毛利、UPC、门店和周次。

因此，本项目不是从零构造一个 toy dataset，而是站在 scanner-data retail pricing 的经典经验研究语境上。README 中应明确写出：DFF 是历史数据，适合方法演示和 portfolio 复现；上线到真实业务时，需要替换为当前业务的 transaction / inventory / campaign / cost 数据。

\subsection{Differentiated Demand 与 Cereal 文献}

BLP (Berry, Levinsohn, and Pakes, 1995) 和 Nevo (2001) 是差异化产品需求估计的核心参照。它们关心的是结构需求、产品替代、市场势力和价格-成本边际，而不是简单预测下周销量。Cereals 品类有多个大品牌、明显产品差异和可解释的交叉替代关系，所以适合把 own-price elasticity 与 cross-price elasticity 作为模型输出。

本项目有意识地不做完整 BLP 或品牌对品牌的结构替代矩阵。MVP 使用 SKU-store FE、week FE、聚合竞争价格和二元促销项，目标是做透明决策支持，而不是估计一个可发表的结构产业组织模型。这个边界需要写清楚，反而会让项目显得更可信。

\subsection{Promotion、库存行为与静态模型边界}

促销不只是价格下降。Hendel and Nevo (2006) 强调，temporary sales 可能引发消费者 stockpiling；如果忽略这种动态行为，静态需求模型可能错误衡量长期价格反应。DFF 的 \texttt{SALE} 字段让本项目可以把促销状态显式纳入模型，但 MVP 仍然是静态周级模型。

因此，文档和 app 必须保留两层 caveat：第一，促销状态本身可能内生；第二，促销周销量提升可能包含跨期购买转移。MVP 的推荐适合做短期决策模拟，不能直接解释为长期均衡销量提升。

\subsection{Revenue Optimization 与商业决策系统}

Revenue management / pricing optimization 文献把需求曲线、边际成本、容量或库存约束放进同一个优化问题。Talluri and van Ryzin (2004) 与 Phillips (2005) 是这条线的经典参考。本项目沿用这个工程逻辑：先估计 price-response，再模拟 counterfactual，再在 margin、inventory、price band 和 promo cost 约束下优化。

与商业定价软件相比，本项目的重点不是覆盖所有企业级功能，而是把黑箱系统拆成一个可读的求职 demo：数据清洗、弹性估计、利润函数、优化器、不确定性、识别 caveat 和 A/B test 验证全部在一个 repo 里闭环。

\subsection{本项目的可讲创新点}

\begin{decisionbox}{定位句}
This is not a new econometric estimator. It is an applied decision product that translates classic scanner-data demand estimation and pricing optimization into a transparent, interactive, experiment-aware tool.
\end{decisionbox}

面试或 README 中可以把差异化讲成四点：

\begin{enumerate}
  \item \textbf{从研究到决策。}不是只报告弹性，而是把弹性放进利润函数和约束优化器；
  \item \textbf{从单品价格到品类竞争。}不是只估 own-price response，而是加入 cereal 品牌间的聚合 cross-price term；
  \item \textbf{从观察性估计到实验验证。}不把 FE 回归包装成最终因果结论，而是输出 A/B test 或 rollout 设计；
  \item \textbf{从 notebook 到 app。}把模型变成业务方可以调 price、promo、cost、inventory 和 competitor price 的交互工具。
\end{enumerate}

\section{公式来源与解释边界}

\subsection{三类公式}

本项目中的公式分成三类，必须在 methodology 和面试叙事中分清：

\begin{enumerate}
  \item \textbf{DFF accounting transformations：}有效价格、销售额和 unit-cost proxy 来自 DFF 字段与 manual 的会计定义；
  \item \textbf{标准计量与定价公式：}log-log fixed-effects demand、Duan smearing、A/B test 样本量、constrained pricing optimization 都来自成熟方法；
  \item \textbf{项目实现选择：}竞争价格指数、cell-anchored counterfactual、grid search 和 scenario overlay 是本项目为决策支持做的工程化实现，不是新的理论模型。
\end{enumerate}

\subsection{DFF 字段到模型变量}

DFF 有时记录 bundle price，因此有效单位价格定义为：

\begin{equation}
P^{eff}_{ist}=\frac{PRICE_{ist}}{QTY_{ist}}.
\end{equation}

对应销售额为：

\begin{equation}
Sales_{ist}=\frac{PRICE_{ist}\times MOVE_{ist}}{QTY_{ist}}.
\end{equation}

unit cost 不是观测到的边际成本，而是由 DFF gross-margin 字段反推的 average-acquisition-cost proxy：

\begin{equation}
c_{ist}=P^{eff}_{ist}\left(1-\frac{PROFIT_{ist}}{100}\right).
\end{equation}

因此，利润模拟是方向性 decision simulation，不应被解释为精确会计利润或 observed marginal-cost optimization。

\subsection{需求、重转换与反事实}

主需求方程是 reduced-form constant-elasticity panel approximation，不是 BLP / Nevo 结构估计，也不是品牌对品牌的完整替代矩阵。BLP 和 Nevo 的作用是提供 differentiated-products demand 与 cereal market 的研究语境，不是本项目方程的直接来源。

Duan smearing correction

\begin{equation}
\widehat{S}=\frac{1}{N}\sum_{n=1}^{N}\exp(\widehat{\epsilon}_n)
\end{equation}

只用于把 fitted log-demand 转回 level prediction。app 优化器中的 anchored counterfactual 从 observed level baseline 出发：

\begin{equation}
\widehat{Q}(p',m',p^{comp\prime})=
\bar{Q}_{cell}\,
\exp\left[
\widehat{\beta}_{own}\Delta\log p
+\widehat{\beta}_{cross}\Delta\log p^{comp}
+\widehat{\theta}\Delta m
\right],
\end{equation}

因此不再额外乘 smearing factor。默认情况下竞争价格固定，$\Delta\log p^{comp}=0$；只有当用户显式运行 competitor-price scenario 时，交叉价格项才进入反事实变化。

\subsection{收入、利润与 guardrail 解释}

inventory cap 必须同时约束 revenue 和 profit。令

\begin{equation}
Q^{sold}(p',m',p^{comp\prime})=
\min\{\widehat{Q}(p',m',p^{comp\prime}),Inventory\},
\end{equation}

则

\begin{align}
Revenue(p',m',p^{comp\prime}) &= p'\,Q^{sold}(p',m',p^{comp\prime}),\\
Profit(p',m',p^{comp\prime}) &= (p'-c)\,Q^{sold}(p',m',p^{comp\prime})-F m'.
\end{align}

当优化器频繁绑定上界 price guardrail 时，输出应解释为 \textbf{raise-and-test candidate generated by a constrained local demand model}，不是 automatic deployment recommendation。guardrail binding 表明在估计的 constant-elasticity 需求曲线和 DFF cost proxy 下，无约束利润目标偏好比历史支持范围更高的价格。

\section{需求模型}

\subsection{透明基准模型}

主模型采用 SKU-store 固定效应和周固定效应的对数需求模型。MVP 中，价格变量是有效单位成交价，促销变量是二元 indicator；连续折扣深度不在同一模型里混用。由于 $\log(Q)$ 要求正销量，主需求模型过滤 $Q>0$ 后估计的是 \textbf{observed positive-sales cells 上的条件需求}，而不是所有可能 UPC-store-week 的无条件品类需求。零销量和缺失 movement week 可在 robustness 中用补全面板、$\log(Q+1)$ 或 availability indicator 处理。

\begin{equation}
\log(Q_{ist})=\alpha_{is} + \gamma_t + \beta_{own}\log(P^{eff}_{ist}) + \beta_{cross}\log(P^{comp}_{ist}) + \theta Promo_{ist} + X_{ist}'\delta + \epsilon_{ist}.
\end{equation}

其中：
\begin{itemize}
  \item $Q_{ist}$ 是商品 $i$ 在门店 $s$、周 $t$ 的销量；
  \item $P^{eff}_{ist}$ 是有效单位销售价格；
  \item $P^{comp}_{ist}$ 是同一 store-week 中竞争品牌的加权平均 \texttt{price\_per\_oz}；
  \item $\alpha_{is}$ 是 SKU-store pair 固定效应，吸收稳定的商品-门店差异；
  \item $\gamma_t$ 是周固定效应，吸收全市场季节性与共同冲击；
  \item $Promo_{ist}$ 是二元促销状态；
  \item $X_{ist}$ 是可选控制项，MVP 可设为空；增强版可加入节假日、品牌 $\times$ 月份季节项、store-specific trend 或滞后需求控制；
  \item $\beta_{own}$ 是自身价格弹性；
  \item $\beta_{cross}$ 是交叉价格弹性；
  \item $\theta$ 是条件 sale-code effect，以 log points 计量；对应百分比效应是 $\exp(\theta)-1$，不是 $\theta$ 本身。
\end{itemize}

\begin{decisionbox}{为什么使用 SKU-store FE}
Dominick's 是 UPC $\times$ store $\times$ week 面板。不同门店的商圈收入、客流、货架陈列和基准价格水平可能长期不同；只用 SKU FE 会把这些门店间价格差异误记为需求弹性。$\alpha_{is}$ 比单独加入 SKU FE 和 store FE 更保守，因为它允许每个商品在每家门店有自己的基准需求水平。
\end{decisionbox}

\subsection{输出解释}

\begin{decisionbox}{弹性解释}
如果 $\widehat{\beta}_{own}=-1.8$，则自身价格上升 1\% 与销量下降约 1.8\% 相关，条件是控制了 SKU-store 固定效应、周固定效应、竞争价格和促销状态。若 $\widehat{\beta}_{cross}>0$，则竞争品牌价格上升与本商品销量上升相关，符合替代品直觉。若 $\widehat{\theta}=0.43$，它表示促销代码对应约 $e^{0.43}-1\approx54\%$ 的条件性销量差异，而不是 43\% 的直接 uplift。这些估计可以支持决策模拟，但不能自动解释为最终因果效应。
\end{decisionbox}

\subsection{重转换修正}

由于模型估计的是 $\log(Q)$，不能直接把线性预测值取 $\exp(\cdot)$ 当作无偏销量预测。MVP 使用 Duan smearing correction：

\begin{equation}
\widehat{S}=\frac{1}{N}\sum_{n=1}^{N}\exp(\widehat{\epsilon}_n).
\end{equation}

给定线性预测值 $\widehat{\eta}$，销量预测为：

\begin{equation}
\widehat{Q}=\widehat{S}\cdot\exp(\widehat{\eta}).
\end{equation}

增强版可以按 promo 状态、价格分位、品类、品牌或 SKU 分组计算 smearing factor，例如分别估计 $\widehat{S}_{m=0}$ 和 $\widehat{S}_{m=1}$，以处理促销周误差方差更大的情况。第一版先用全样本修正，避免系统性低估销量、收入和利润。

\subsection{识别 caveat}

价格和促销通常都不是随机设定。高需求时期公司可能涨价，低需求时期可能促销；促销也可能来自库存压力、厂家补贴、category captain 排期、竞争对手行动或门店层面的需求预测。这会导致价格和促销系数都存在内生性。因此，文档和 app 必须写清楚：

\begin{quote}
Observed price and promotion variation are not automatically causal. Own-price elasticity, cross-price elasticity, and conditional sale-code effects are used for decision support and should be validated through randomized price/promotion tests, staggered rollouts, credible instruments, or difference-in-differences designs before full rollout.
\end{quote}

\subsection{识别策略升级}

MVP 先使用 SKU-store FE、week FE、竞争价格和促销控制，重点交付透明的决策闭环。增强版把识别问题单独作为模块，而不是把观察性回归包装成最终因果结论。

\begin{tabularx}{\textwidth}{L{3.5cm}Y}
\toprule
策略 & 实现方式与作用 \\
\midrule
价格残差诊断 & 回归 $\log(P^{eff}_{ist})$ on promo、SKU-store FE、week FE，报告 first-stage-like $F$ statistic 和 $R^2$，说明价格变动有多少来自促销与固定效应之外 \\
Hausman IV & 用同一 UPC 或品牌在其他门店同周的价格作为 focal store 价格工具变量；适合需求冲击主要是 store-specific 的设定，但需要警惕全城促销或共同需求冲击 \\
成本冲击 IV & 使用滞后 unit cost 或外部谷物 / 包装材料 PPI 作为工具变量；不直接使用 contemporaneous unit cost，因为它由 price 和 margin 机械反推，可能与价格共同决定 \\
实验验证 & 对候选价格或促销状态做小流量随机实验、geo rollout 或 staggered rollout，把观察性建议转化为可验证决策 \\
\bottomrule
\end{tabularx}

\section{反事实模拟器}

\subsection{预测需求}

给定新的有效单位价格 $p'$、促销状态 $m'\in\{0,1\}$ 和竞争价格 $p^{comp}$，模型层面的 level 预测为：

\begin{equation}
\widehat{Q}(p',m',p^{comp})=\widehat{S}\cdot\exp\left(\widehat{\alpha}_{is} + \widehat{\gamma}_t + \widehat{\beta}_{own}\log(p') + \widehat{\beta}_{cross}\log(p^{comp}) + \widehat{\theta}m' + X_{ist}'\widehat{\delta}\right).
\end{equation}

实际 app 优化器使用 cell-anchored relative response，先从该 cell 的历史平均销量出发，只应用候选价格、竞争价格和促销状态相对于 baseline 的变化：

\begin{equation}
\widehat{Q}_{anchor}(p',m',p^{comp\prime})=
\bar{Q}_{cell}\exp\left[
\widehat{\beta}_{own}\Delta\log p
+\widehat{\beta}_{cross}\Delta\log p^{comp}
+\widehat{\theta}\Delta m
\right].
\end{equation}

默认 MVP 中竞争价格保持不变，即 $\Delta\log p^{comp}=0$；只有当用户显式运行 competitor-price scenario 时，交叉价格项才进入反事实变化。因为 $\bar{Q}_{cell}$ 已经是 level quantity，anchored 反事实不再额外乘 smearing factor，避免把 level correction 重复计算。

MVP 不把连续折扣 $d$ 同时放进 $\log(p'(1-d))$ 和促销项，避免重复计算折扣效应。若增强版需要 discount slider，必须先构造并估计 \texttt{discount\_depth} 的连续系数，再用对应模型做反事实。竞争价格 $p^{comp}$ 默认取当前 store-week 的竞争品牌加权 \texttt{price\_per\_oz}；高级情景可以让用户模拟竞争对手价格变化。

\subsection{收入和利润}

\begin{align}
Q^{sold}(p',m',p^{comp}) &= \min\{\widehat{Q}_{anchor}(p',m',p^{comp}), Inventory\},\\
Revenue(p',m',p^{comp}) &= p'\cdot Q^{sold}(p',m',p^{comp}), \\
Profit(p',m',p^{comp}) &= \left[p'-c\right]\cdot Q^{sold}(p',m',p^{comp})-F\cdot m'.
\end{align}

其中 $c$ 是单位成本。DFF 中 $c$ 由有效单位价格和 gross margin 反推，是 average-acquisition-cost proxy，不是 observed marginal cost。库存字段 DFF 不提供，因此 $Inventory$ 作为 scenario input。$F$ 是可选促销固定成本，MVP 默认 $F=0$，增强版可由用户输入陈列、人工或 coupon 运营成本。注意 revenue 和 profit 都必须使用同一个 capped sold quantity，不能 revenue 用 predicted demand、profit 才 cap inventory。

当 app 展示候选动作的 absolute profit 时，促销固定成本直接进入候选利润：

\begin{equation}
Profit(p',m')=(p'-c)Q^{sold}(p',m')-Fm'.
\end{equation}

当 app 展示相对于观察到的 baseline $(p_0,m_0)$ 的 profit lift 时，促销固定成本必须差分：

\begin{equation}
\Delta Profit =
\left[(p'-c)Q^{sold}(p',m')-Fm'\right]
-
\left[(p_0-c)Q^{sold}(p_0,m_0)-Fm_0\right].
\end{equation}

因此，$F$ 只有在候选动作和 baseline 的促销状态不同的时候才改变 incremental lift；如果二者都促销或都不促销，固定成本在差分中抵消。

\subsection{Scenario overlay 与 shock 定义域}

Streamlit app 在冻结的需求模型之上提供 business shock overlay，而不是重新估计模型。记 $\delta_q$ 为需求冲击、$\delta_c$ 为成本冲击、$\delta_{comp}$ 为竞争价格冲击、$\bar{Q}$ 为可选库存上限，则：

\begin{align}
\Delta\log p^{comp} &= \log(1+\delta_{comp}),\\
\widetilde{Q}(p',m') &= \widehat{Q}_{anchor}\!\left(p',m';\,\Delta\log p^{comp}\right)(1+\delta_q),\\
Q^{sold}(p',m') &= \min\{\widetilde{Q}(p',m'),\bar{Q}\},\\
c^{eff} &= c(1+\delta_c),\\
Revenue(p',m') &= p'Q^{sold}(p',m'),\\
Profit(p',m') &= (p'-c^{eff})Q^{sold}(p',m')-Fm'.
\end{align}

这些 multiplicative shocks 必须满足：

\begin{equation}
\delta_q>-1,\qquad \delta_c>-1,\qquad \delta_{comp}>-1,
\end{equation}

否则需求、成本或竞争价格会变成负数，且 $\Delta\log p^{comp}$ 不再有定义。app 的 risk flags 使用对称阈值：

\begin{equation}
|\delta_q|>0.20,\quad |\delta_c|>0.25,\quad
|\delta_{comp}|>0.15,\quad \text{or}\quad
\bar Q < \bar Q^{base}_{cell}.
\end{equation}

\subsection{输出表}

\begin{tabularx}{\textwidth}{L{3.0cm}C{2.0cm}C{2.0cm}C{2.4cm}C{2.3cm}C{2.3cm}}
\toprule
情景 & 有效价格 & 促销状态 & 预测销量 & 收入 & 利润 \\
\midrule
当前策略 & $p_0$ & $m_0$ & $Q_0$ & $R_0$ & $\Pi_0$ \\
用户输入策略 & $p'$ & $m'$ & $\widehat{Q}$ & $\widehat{R}$ & $\widehat{\Pi}$ \\
最优策略 & $p^*$ & $m^*$ & $\widehat{Q}^*$ & $\widehat{R}^*$ & $\widehat{\Pi}^*$ \\
\bottomrule
\end{tabularx}

\section{优化器设计}

\subsection{目标函数}

app 提供三个目标：

\begin{enumerate}
  \item \textbf{最大化收入：}
  \begin{equation}
  \max_{p,m} Revenue(p,m,p^{comp}).
  \end{equation}
  \item \textbf{最大化利润：}
  \begin{equation}
  \max_{p,m} Profit(p,m,p^{comp}).
  \end{equation}
  \item \textbf{最大化利润且满足销量约束：}
  \begin{equation}
  \max_{p,m} Profit(p,m,p^{comp}) \quad \text{s.t. } Q^{sold}(p,m,p^{comp})\ge Q_{\min}.
  \end{equation}
\end{enumerate}

\subsection{业务约束}

\begin{align}
p_{\min} &\le p \le p_{\max},\qquad
m\in\{0,1\},\qquad
p-c \ge margin_{\min},\\
Q^{sold}(p,m,p^{comp})
&= \min\{\widehat{Q}_{anchor}(p,m,p^{comp}),Inventory\}.
\end{align}

约束必须存在，否则优化器会像玩具。真实业务中不能推荐负毛利、过低有效价格、超过库存承载能力或破坏品类价格带的策略。若预测需求超过库存，app 使用 $Q^{sold}$ 计算可实现收入和利润，并把 stockout risk 作为 guardrail flag。若增强版加入连续折扣深度，再额外加入 $0 \le d \le d_{\max}$。

\subsection{实现方式}

MVP 用 grid search。因为 $m$ 是二元变量，优化实际是两次一维搜索：先固定 $m=0$ 扫描价格，再固定 $m=1$ 扫描价格，最后比较两个最优利润。理由：透明、稳定、容易解释，面试中也能清楚说明复杂度。

增强版可用 \texttt{scipy.optimize}，但不是必要条件。先把业务闭环做完整，再提高算法复杂度。

\begin{decisionbox}{上界绑定是诊断，不是部署建议}
04 counterfactual 的核心警示是：98.5\% 的 eligible cells 推荐价格贴在上界 guardrail。这不是“优化器很自信”的证据，而是 constant-elasticity demand 与当前 cost proxy 共同推出的诊断结果。对于 interior single-product optimum，若边际成本恒定且 $\varepsilon<-1$，constant-elasticity demand $Q=A p^\varepsilon$ 推出 Lerner condition：
\[
\frac{p-c}{p}=-\frac{1}{\varepsilon}.
\]
当 $\varepsilon\approx -1.73$ 时，隐含无约束 margin 约为 58\%，因此模型会机械地偏好高价。正确 framing 是 raise-and-test candidates，不是 production deployment recommendations。
\end{decisionbox}

\section{不确定性与稳健性}

\subsection{必须展示不确定性}

候选价格附近必须展示利润置信区间或敏感性区间。可选实现：

\begin{itemize}
  \item 对需求模型参数做 bootstrap；
  \item 用弹性估计的置信区间模拟上下界；
  \item 对成本和库存做 scenario sensitivity；
  \item 比较不同模型设定下的候选价格是否稳定。
\end{itemize}

\subsection{稳健性检查}

\begin{tabularx}{\textwidth}{L{4.0cm}Y}
\toprule
检查项 & 目的 \\
\midrule
去除极端价格周 & 检查促销或清仓是否驱动弹性 \\
比较 SKU FE、store FE、SKU-store FE & 检查门店间价格差异是否污染自身价格弹性 \\
加入 week/month fixed effects & 控制季节性与全市场共同冲击 \\
加入 promo 控制 & 避免把条件 sale-code effect 误记为价格效应；同时承认 promo 仍可能内生 \\
加入竞争价格 & 检查 Cereals 品牌替代效应是否影响候选价格 \\
交叉弹性符号检查 & 若 $\widehat{\beta}_{cross}\le 0$，标记为不稳定估计；优化器使用敏感性区间或关闭交叉价格响应，而不是机械代入可疑系数 \\
用 holdout 验证预测误差 & 防止模型只在样本内表现好 \\
弹性上下界模拟 & 检查候选价格是否对参数敏感 \\
\bottomrule
\end{tabularx}

额外建议加入 store-week fixed effects 规格，作为比 week FE 更严格的 confounding 检查：

\begin{equation}
\log(Q_{ist}) =
\alpha_{is}
+
\lambda_{st}
+
\beta_{own}\log(P^{eff}_{ist})
+
\theta Promo_{ist}
+
\epsilon_{ist}.
\end{equation}

其中 $\lambda_{st}$ 吸收本地需求冲击、门店层面的 campaign、特定周的库存或客流条件。这个规格很重要，因为 DFF 的 price / promotion 是 manager choice，且 \texttt{SALE} code 不是干净的 randomized treatment；manual 也提示，即使 sale code 没有设置，当周仍可能存在促销。store-week FE 会牺牲一部分 cross-price variation，但可以更严格检查 own-price elasticity 是否由 local store-week confounding 驱动。若竞争价格指数在该规格中被本地品类价格变动机械吸收、共线或不稳定，则 store-week-FE robustness 中应排除 $P^{comp}_{ist}$。

\section{A/B test 验证模块}

\subsection{为什么必须有实验模块}

观察性需求估计不能替代线上实验。候选价格上线前应通过小流量实验验证。实验模块让该 demo 从“模型练习”变成“可落地决策系统”。

\subsection{实验输出}

\begin{itemize}
  \item treatment：推荐有效价格 / 促销状态；
  \item control：当前有效价格 / 促销状态；
  \item primary metric：profit per visitor 或 revenue per visitor；
  \item guardrail metrics：conversion、return rate、stockout、customer complaints；
  \item sample size；
  \item test duration；
  \item randomization unit：user、session、geo、SKU-cluster。
\end{itemize}

\subsection{样本量公式}

\begin{equation}
n = \frac{2(z_{1-\alpha/2}+z_{1-\beta})^2\sigma^2}{\Delta^2}.
\end{equation}

其中 $\alpha$ 是显著性水平，$1-\beta$ 是 power，$\Delta$ 是最小可检测效应（MDE），$\sigma^2$ 是主要指标方差。

\section{Web app 设计}

\subsection{技术选择}

\begin{decisionbox}{技术栈}
使用 Python + Streamlit。不要一开始做 React 全栈。求职信号来自模型闭环和业务解释，不来自前端复杂度。
\end{decisionbox}

\begin{tabularx}{\textwidth}{L{3.2cm}Y}
\toprule
层级 & 工具 \\
\midrule
数据处理 & pandas, numpy \\
统计建模 & statsmodels, linearmodels（可选） \\
机器学习 benchmark & scikit-learn（可选） \\
优化 & scipy.optimize 或 grid search \\
可视化 & plotly 或 matplotlib \\
Web app & Streamlit \\
部署 & Streamlit Community Cloud / Render \\
版本管理 & GitHub \\
文档 & README + methodology doc + model card + executive memo \\
\bottomrule
\end{tabularx}

\subsection{页面结构}

\begin{center}
\small
\begin{tabularx}{\textwidth}{L{3.4cm}L{4.3cm}Y}
\toprule
页面 & 用户操作 & 输出 \\
\midrule
Demand Model & 查看模型结果 & 自身弹性、交叉弹性、sale-code log-point effect、holdout error、识别 caveat \\
Counterfactual Simulator & 调整 effective price、promo、competitor price、cost、inventory & capped sold units、收入、利润、变化率、风险提示 \\
Profit Optimizer & 点击 rerun optimizer 或调整 scenario & $p^*$、$m^*$、$\Delta Revenue$、$\Delta Profit$、upper-guardrail flag \\
Experiment Design & 查看候选实验或输入 MDE / power & 样本量、实验周期、guardrail metrics、underpowered flag \\
Limitations & 阅读模型限制 & 价格/促销内生性、AAC cost proxy、库存场景化、raise-and-test framing \\
Upload \& Score & 上传业务方 CSV & 字段校验、scenario scoring、候选动作和风险提示 \\
\bottomrule
\end{tabularx}
\end{center}

\subsection{关键交互}

\begin{itemize}
  \item slider：new effective unit price；
  \item slider：competitor price per oz；
  \item toggle：promo on / off；
  \item number input：unit cost；
  \item number input：promo fixed cost；
  \item number input：inventory；
  \item selectbox：objective（revenue / profit / constrained profit）；
  \item button：Run Counterfactual；
  \item button：Find Optimal Price / Re-run optimizer；
  \item file uploader：上传自有 SKU 数据并映射到 app schema；
  \item warning box：observational estimates require experimental validation。
\end{itemize}

\section{关键图表}

\begin{tabularx}{\textwidth}{L{4.0cm}Y}
\toprule
图表 & 作用 \\
\midrule
Historical price and quantity & 展示 panel 数据结构和时间变化 \\
Price vs quantity scatter & 直观展示相关性，但注明不是因果 \\
Estimated demand curve & 展示价格到需求的函数关系 \\
Revenue curve & 显示收入最大化价格 \\
Profit curve & 显示利润最大化价格；这是核心图 \\
Baseline vs raise-and-test candidate & 把模型结果翻译成业务指标，同时标注部署风险 \\
Uncertainty band around profit curve & 展示模型风险和决策稳健性 \\
Sample size curve & 展示 A/B test 所需样本随 MDE 变化 \\
\bottomrule
\end{tabularx}

\section{项目文件结构}

\begin{lstlisting}
pricing-decision-engine/
|
|-- app.py
|-- README.md
|-- requirements.txt
|-- pricing_promotion_decision_engine_plan.tex
|-- pricing_promotion_decision_engine_plan.pdf
|
|-- docs/
|   |-- app_wireframe.md
|   |-- methodology.md
|   |-- methodology.tex
|   |-- methodology.pdf
|
|-- rawData/
|   |-- README.md
|   |-- dominicks_manual.pdf
|   |-- wcer.csv
|   |-- upccer.csv
|
|-- data/
|   |-- processed/
|       |-- sku_store_week_panel.parquet
|       |-- brand_size_store_week_panel.parquet
|       |-- demand_modeling_dataset.parquet
|       |-- model_coefficients.csv
|       |-- top_recommendations.csv
|       |-- experiment_candidates.csv
|
|-- notebooks/
|   |-- 01_data_cleaning.ipynb
|   |-- 02_eda.ipynb
|   |-- 03_demand_estimation.ipynb
|   |-- 04_counterfactual.ipynb
|   |-- 05_ab_testing_design.ipynb
|
|-- pages/
|   |-- 1_Demand_Model.py
|   |-- 2_Counterfactual_Simulator.py
|   |-- 3_Profit_Optimizer.py
|   |-- 4_Experiment_Design.py
|   |-- 5_Limitations.py
|   |-- 6_Upload_and_Score.py
|
|-- src/
|   |-- data.py
|   |-- features.py
|   |-- validation.py
|   |-- scenario.py
|   |-- simulation.py
|   |-- optimization.py
|   |-- plots.py
|   |-- upload.py
|
|-- tests/
|   |-- test_features.py
|   |-- test_scenario.py
|   |-- test_optimization.py
|   |-- test_upload.py
|
|-- reports/
    |-- data_cleaning_summary.md
    |-- eda_summary.md
    |-- demand_model_summary.md
    |-- counterfactual_summary.md
    |-- ab_test_plan.md
    |-- figures/
\end{lstlisting}

\section{README 结构}

README 不应写成技术流水账，而应写成商业 memo + 技术说明。

根目录 \texttt{README.md} 是整个项目的首页；\texttt{rawData/README.md} 只写原始数据来源、文件清单、字段字典、manual 摘要和读取规则；\texttt{reports/data\_cleaning\_summary.md} 记录真实清洗诊断与样本规模。不要把根 README 写成 codebook，也不要把 raw data README 写成项目 pitch。

\subsection{必须包含的部分}

\begin{enumerate}
  \item \textbf{Business Problem：}定价和促销如何影响需求、收入、利润和库存风险。
  \item \textbf{Data：}数据来源、粒度、商品数量、观察值数量、字段说明、真实字段与 scenario inputs。
  \item \textbf{Relation to Prior Work：}说明项目与 DFF scanner-data 研究、cereal demand 文献、revenue optimization 和 experimentation 的关系；强调这是 applied decision product，不是新结构估计器。
  \item \textbf{Method：}需求估计、反事实模拟、约束优化、A/B test 设计。
  \item \textbf{Results：}弹性、raise-and-test 候选价格、收入变化、利润变化、样本量。
  \item \textbf{How to Run：}安装依赖和启动 app。
  \item \textbf{Limitations：}价格和促销内生性、竞争价格构造假设、库存场景化、上线前需要实验验证。
\end{enumerate}

\subsection{README 必须写明的模型限制}

\begin{quote}
The aggregate cross-price term approximates average substitution intensity across competing cereal brands. A full brand-pair substitution matrix would require a discrete-choice or random-coefficients structural demand model and is outside the MVP scope.
\end{quote}

\begin{quote}
This app is a transparent decision-support implementation inspired by scanner-data demand estimation and pricing optimization. It should not be presented as a new econometric estimator or as a replication of Nevo's structural cereal model.
\end{quote}

\begin{quote}
The competitor price index uses baseline or lagged sales-share weights rather than current-week sales shares, reducing mechanical endogeneity from contemporaneous demand shocks. Results should still be interpreted as decision support, not final causal evidence.
\end{quote}

\subsection{README 结果表模板}

\begin{tabularx}{\textwidth}{L{5.0cm}Y}
\toprule
Result & Value \\
\midrule
Estimated own-price elasticity & 来自模型，不手填 \\
Estimated cross-price elasticity & 来自模型，不手填 \\
Promo sale-code effect & 来自模型，不手填；若以百分比展示，使用 $\exp(\theta)-1$ \\
Raise-and-test candidate price & 来自优化器 \\
Expected revenue change & 来自反事实模拟 \\
Expected profit change & 来自反事实模拟 \\
Required sample size for A/B test & 来自实验模块 \\
\bottomrule
\end{tabularx}

\newpage
\section{执行路线图}

\small
\begin{longtable}{C{1.1cm}L{2.8cm}L{7.0cm}L{3.2cm}}
\toprule
阶段 & 任务 & 具体动作 & 产出 \\
\midrule
\endfirsthead
\toprule
阶段 & 任务 & 具体动作 & 产出 \\
\midrule
\endhead
1 & 数据选择与清洗 & 已完成 DFF Cereals 清洗、SALE code 处理、unit price / cost proxy 构造 & SKU-store panel \\
2 & EDA 与聚合 & 已完成 brand mapping、size parsing、brand-size panel、粒度选择诊断 & \texttt{02\_eda.ipynb} + EDA reports \\
3 & 需求估计 & 已估计 log demand with SKU-store FE, week FE, own price, competitor price, promo & coefficients + holdout metrics \\
4 & 反事实函数 & 已实现 anchored counterfactual，默认固定 competitor price，可接受 scenario shock & \texttt{simulation.py} + \texttt{scenario.py} \\
5 & 优化器 & 已实现 constrained grid search，含 price guardrail、margin、inventory、promo cost & \texttt{optimization.py} \\
6 & 实验模块 & 已生成 A/B validation design、sample size curve、risk flags & 05 notebook \\
7 & Streamlit app & 已搭建六页 MVP，含 simulator、optimizer、limitations、upload scoring & local Streamlit app \\
8 & 文档同步 & 已同步 README、methodology、reports 和本计划的公式口径 & GitHub-ready docs \\
9 & 下一步 A & 检查同品牌跨包装规格 cannibalization，判断 98.5\% guardrail binding 是否稳健 & 07 notebook \\
10 & 下一步 B/C & 补 tests、部署、case study；可选 Hausman IV / cost-shock IV & portfolio link + causal-IV extension \\
\bottomrule
\end{longtable}
\normalsize

\section{一页 Executive Memo 模板}

\begin{actionbox}{Memo 结构}
\textbf{Problem.} Pricing and promotion decisions affect demand, revenue, margin, and inventory risk. The current analysis builds a decision engine to evaluate pricing counterfactuals and generate constrained raise-and-test pricing candidates.\\[0.35em]
\textbf{Relation to prior work.} The project is positioned as an applied decision product built on scanner-data demand estimation, cereal differentiation research, revenue optimization, and online experimentation, not as a new econometric estimator.\\[0.35em]
\textbf{Method.} Estimate demand response using SKU-store-week panel data with SKU-store and week fixed effects, own price, competitor price, and promotion controls; use Duan smearing only for log-demand level prediction; use cell-anchored relative response for counterfactual optimization; compute revenue and profit from inventory-capped sold units; design an A/B test for validation.\\[0.35em]
\textbf{Candidate action.} Report the raise-and-test effective price and promotion status, expected revenue/profit change, upper-guardrail flag, and decision risk. Values must be generated by the app.\\[0.35em]
\textbf{Caveat.} Observational price and promotion variation is not final causal evidence. Candidate actions should be validated with randomized or quasi-experimental rollout.
\end{actionbox}

\section{简历 bullet 与面试话术}

\subsection{简历 bullet}

\begin{quote}
Built an interactive pricing and promotion decision engine using Python, Streamlit, and SKU-store-week panel demand estimation; estimated own- and cross-price elasticity, simulated effective price and promotion counterfactuals, generated constrained raise-and-test pricing candidates under cost and inventory constraints, and designed an A/B testing framework for validation.
\end{quote}

如果已有真实计算结果，可以改为：

\begin{quote}
Built an interactive pricing and promotion decision engine using Python, Streamlit, and SKU-store-week panel data; estimated own- and cross-price elasticity, flagged upper-guardrail optimizer behavior, and quantified expected revenue and margin changes under alternative pricing and promotion scenarios.
\end{quote}

不要写无法证明的百分比。百分比必须来自 app 的实际输出。

\subsection{60 秒面试版本}

\begin{quote}
I built a pricing decision engine for retail products using Dominick's SKU-store-week scanner data. It is positioned as an applied decision product, not a new econometric estimator: the app takes ideas from scanner-data demand estimation, cereal differentiation research, and revenue optimization, then estimates own-price elasticity, cross-price elasticity, and promotion response with fixed effects. It simulates counterfactual pricing and promotion strategies, generates constrained raise-and-test candidates under cost and inventory constraints, and includes an A/B testing module because observational price and promotion variation is not enough for causal rollout decisions.
\end{quote}

\section{关键风险清单}

\begin{riskbox}{项目失败的主要原因}
\begin{enumerate}
  \item \textbf{变成销量预测项目。}预测只是输入，不是最终产品。
  \item \textbf{没有利润函数。}只看 revenue 会被定价团队质疑。
  \item \textbf{忽略价格和促销内生性。}必须写 caveat、识别诊断和实验验证方案。
  \item \textbf{没有约束。}真实定价有库存、成本、价格下限、最低 margin。
  \item \textbf{没有 web app。}notebook 不足以证明决策工具化能力。
  \item \textbf{没有 README 和 memo。}招聘方不会自己读代码推断你的商业价值。
\end{enumerate}
\end{riskbox}

\section{最终交付物清单}

\begin{tabularx}{\textwidth}{L{4.0cm}Y L{2.4cm}}
\toprule
交付物 & 内容 & 优先级 \\
\midrule
Streamlit web app & 交互式定价、促销、反事实、优化、实验设计 & 必须 \\
GitHub repo & 可运行代码、清楚结构、requirements、核心函数 tests & 必须 \\
README & 商业问题、数据、方法、结果、限制、运行方式 & 必须 \\
Methodology doc & 需求模型、反事实模拟、利润函数、优化器、识别 caveat，提供 Markdown 与 TeX 两版 & 必须 \\
Executive memo & 一页业务结论，适合发给内推或 recruiter & 必须 \\
Limitations / model card 页面 & 数据限制、识别限制、适用范围、guardrail binding 风险 & 建议 \\
Demo video/GIF & 60-90 秒演示 app 如何使用 & 建议 \\
Resume bullet & 对齐定价、因果、优化、实验 & 必须 \\
\bottomrule
\end{tabularx}

\section{最终判断}

这个 demo 的正确形式是：

\begin{decisionbox}{核心闭环}
\centering
\textbf{数据} $\rightarrow$ \textbf{需求估计} $\rightarrow$ \textbf{反事实模拟} $\rightarrow$ \textbf{利润优化} $\rightarrow$ \textbf{实验验证} $\rightarrow$ \textbf{业务建议}
\end{decisionbox}

最终输出必须让招聘方在 60 秒内看懂：你不是只会做学术结构模型，也不是只会做 dashboard；你能把经济学、因果推断和数据科学压缩成一个直接影响价格、促销和利润的决策工具。

\section*{参考文献}
\addcontentsline{toc}{section}{参考文献}

\begin{enumerate}
  \item Angrist, Joshua D., and Jörn-Steffen Pischke. 2009. \textit{Mostly Harmless Econometrics: An Empiricist's Companion}. Princeton University Press.
  \item Berry, Steven, James Levinsohn, and Ariel Pakes. 1995. ``Automobile Prices in Market Equilibrium.'' \textit{Econometrica} 63(4): 841--890. \url{https://doi.org/10.2307/2171802}.
  \item Duan, Naihua. 1983. ``Smearing Estimate: A Nonparametric Retransformation Method.'' \textit{Journal of the American Statistical Association} 78(383): 605--610. \url{https://doi.org/10.1080/01621459.1983.10478017}.
  \item Hendel, Igal, and Aviv Nevo. 2006. ``Sales and Consumer Inventory.'' \textit{RAND Journal of Economics} 37(3): 543--561. \url{https://doi.org/10.1111/j.1756-2171.2006.tb00030.x}.
  \item Imbens, Guido W., and Donald B. Rubin. 2015. \textit{Causal Inference for Statistics, Social, and Biomedical Sciences}. Cambridge University Press.
  \item Kohavi, Ron, Diane Tang, and Ya Xu. 2020. \textit{Trustworthy Online Controlled Experiments: A Practical Guide to A/B Testing}. Cambridge University Press.
  \item Nevo, Aviv. 2001. ``Measuring Market Power in the Ready-to-Eat Cereal Industry.'' \textit{Econometrica} 69(2): 307--342. \url{https://doi.org/10.1111/1468-0262.00194}.
  \item Phillips, Robert L. 2005. \textit{Pricing and Revenue Optimization}. Stanford University Press.
  \item Talluri, Kalyan T., and Garrett J. van Ryzin. 2004. \textit{The Theory and Practice of Revenue Management}. Springer.
  \item Train, Kenneth E. 2009. \textit{Discrete Choice Methods with Simulation}. Cambridge University Press.
  \item UChicago Booth Kilts Center for Marketing. n.d. ``Dominick's Finer Foods Database.'' Accessed April 2026. \url{https://www.chicagobooth.edu/research/kilts/research-data/dominicks}.
\end{enumerate}

\end{document}
